% Operadores:
\DeclareMathOperator*{\argmin}{arg\,min}
\DeclareMathOperator*{\argmax}{arg\,max}
\DeclareMathOperator*{\essinf}{ess\,inf}

% Insertar imagen:
\newcommand{\insertimage}[3]{
    \begin{figure}
        \centering
        \includegraphics[width=#2\textwidth]{./images/#1}
        \caption{#3}
        \label{fig:#1}
    \end{figure}
}

% Ecuaciones (revisar no utilizados):
\newcommand{\parent}[1]{\left(#1\right)}
\newcommand{\rparent}[1]{\left[#1\right]}
\newcommand{\feasible}[2]{\genfrac {}{}{0pt}{2}{#1}{#2}}
\newcommand{\lspace}[2]{\operatorname{L}^{#1}\parent{#2}}
\newcommand{\borel}[1]{\mathcal{B}\parent{#1}}
\newcommand{\TV}[1]{\norm{#1}_{\operatorname{TV}}}
\newcommand{\bernoulli}[1]{\operatorname{Bernoulli}\parent{#1}}
\newcommand{\E}[2]{\mathbb{E}_{#1}\rparent{#2}}
\newcommand{\var}[1]{\operatorname{Var}\parent{#1}}
\newcommand{\cov}[2]{\operatorname{Cov}\parent{#1,\,#2}}
\newcommand{\ley}[1]{\operatorname{Ley}\parent{#1}}
\newcommand{\KL}[2]{\operatorname{D_{KL}}\parent{#1\,\|\,#2}}
\newcommand{\fdiv}[2]{\operatorname{D}_f\parent{#1\,\|\,#2}}
\newcommand{\tv}[2]{\operatorname{D_{TV}}\parent{#1\,\|\,#2}}
\newcommand{\DF}[2]{\operatorname{D_F}\parent{#1\,\|\,#2}}
\newcommand{\entropy}[1]{\mathcal{H}\parent{#1}}
\newcommand{\elbo}{\operatorname{ELBO}}
\newcommand{\cte}{\operatorname{constante}}
\newcommand{\trace}[1]{\operatorname{Tr}\parent{#1}}
\newcommand{\supp}[1]{\operatorname{Supp}\parent{#1}}
\newcommand{\score}[2]{\nabla_{#1} \log #2}
\newcommand{\R}{\mathbb{R}}
\newcommand{\N}{\mathbb{N}}
\newcommand{\matrixspace}[2]{\mathcal{M}_{#1}\parent{#2}}
\newcommand{\dotproduct}[2]{\left\langle #1,\,#2\right\rangle}
\newcommand{\norm}[1]{\left\lVert #1 \right\rVert}
\newcommand{\weak}{\rightharpoonup}
\newcommand{\gaussian}[2]{\mathcal{N}\parent{#1,\,#2}}
\newcommand{\gaussianexpl}[3]{\mathcal{N}\parent{#1;\,#2,\,#3}}
\newcommand{\identity}[1]{\operatorname{I}_{#1}}
\newcommand{\1}{\mathds{1}}
\newcommand{\diag}[1]{\operatorname{diag}\parent{#1}}
\newcommand{\hessian}[1]{\operatorname{D}^2 #1}
\newcommand{\continuous}[1]{\mathcal{C}\parent{#1}}
\newcommand{\continuousb}[1]{\mathcal{C}_b\parent{#1}}
\newcommand{\posmeasure}[1]{\mathcal{M}_+\parent{#1}}
\newcommand{\discretemeasure}[1]{\mathcal{M}_+^d\parent{#1}}
\newcommand{\probmeasure}[1]{\mathcal{M}_+^1\parent{#1}}
\newcommand{\probmeasurep}[1]{\mathcal{M}_+^{1,p}\parent{#1}}
\newcommand{\xspace}{\mathcal{X}}
\newcommand{\yspace}{\mathcal{Y}}
\newcommand{\salgebra}{\mathcal{S}}
\newcommand{\simplex}[1]{\Sigma_{#1}}
\newcommand{\wasserstein}[2]{\mathcal{W}_{#1}\parent{#2}}
\newcommand{\pendiente}[1]{\colorbox{yellow}{#1}} % temporal.
\newcommand{\monge}[1]{\operatorname{Monge}\parent{#1}}
\newcommand{\kantorovich}[1]{\operatorname{Kantorovich}\parent{#1}}
\newcommand{\dual}[1]{\operatorname{J}\parent{#1}}
\newcommand{\kantoroviche}[1]{\operatorname{Kantorovich}_\epsilon\parent{#1}}
\newcommand{\gibbs}{\mathcal{K}}
\newcommand{\C}{\mathcal{C}}
\newcommand{\power}[1]{\mathcal{P}\parent{#1}}
\newcommand{\ptrue}{p_{\operatorname{true}}}
\newcommand{\pprior}{p_{\operatorname{prior}}}
\renewcommand{\P}{\mathbb{P}}
\renewcommand{\d}{\,\operatorname{d}\!}





